\section{Structural Tools and Regime Tests}
\label{app:structural-tools-regime-tests}
\addcontentsline{toc}{section}{Appendix: Structural Tools and Regime Tests}

This appendix collects the main structural tools extracted from the
proof architecture. Their purpose is not to replace the main
derivations, but to normalize frequently recurring arguments into
reusable tests and reduction principles.

\subsection*{Sector and closure tools}

\paragraph{Tool (Closure-transfer principle).}
Let $X$ be any operator on the septonic probe space. Its admissible
reduction is
\[
  X_{\mathrm{adm}} := \hat{P}_{C_7} X \hat{P}_{C_7}.
\]
For admissible states $\Psi \in \operatorname{Ran}(\hat{P}_{C_7})$, the
physically relevant action of $X$ is determined by $X_{\mathrm{adm}}$,
while the complementary blocks measure off-sector leakage.

\paragraph{Tool (Equivalent closure criteria).}
The following are equivalent
(cf.\ Lemma~\texttt{lem:closure-forcing-lemma} in the active-core PDF):
\begin{enumerate}[label=(\roman*)]
  \item $\mathfrak{R}_C = 0$;
  \item $[\mathcal{M}_{16}, \hat{P}_{C_7}] = 0$;
  \item $\mathcal{M}_{16}$ preserves the admissible sector and its complement.
\end{enumerate}
Hence exact closure is equivalent to block preservation, i.e. block-diagonalization of $\mathcal{M}_{16}$ with respect to
$\mathcal{H}_{\mathrm{sept}} = \hat{P}_{C_7}\mathcal{H}_{\mathrm{sept}}
\oplus (\mathbf{1} - \hat{P}_{C_7})\mathcal{H}_{\mathrm{sept}}$.

\paragraph{Tool (Closure-preservation test).}
For any new structural term $X$, first test whether
$[X, \hat{P}_{C_7}] = 0$.
If yes, then $X$ respects the admissible stabilized sector directly.
If not, then $X$ either produces leakage or must be controlled by a
compensator-mediated defect term
(cf.\ Hypothesis~\texttt{hyp:frame-dependent-defect-readout} in the active-core PDF).

\paragraph{Tool (Sectorial normal form).}
Every operator $X$ decomposes as
\[
  X = \hat{P}_{C_7}X\hat{P}_{C_7}
    + (\mathbf{1}-\hat{P}_{C_7})X\hat{P}_{C_7}
    + \hat{P}_{C_7}X(\mathbf{1}-\hat{P}_{C_7})
    + (\mathbf{1}-\hat{P}_{C_7})X(\mathbf{1}-\hat{P}_{C_7}),
\]
isolating: (i) admissible core block, (ii) incoming leakage,
(iii) outgoing leakage, (iv) complement-sector block.

\subsection*{Master-equation tools}

\paragraph{Tool (Core-defect normal form).}
The projected structural field operator
$\mathfrak{F}_{\mathrm{struct}}
= \hat{P}_{C_7}(i\Gamma^\mu\mathcal{D}_\mu - \Omega)\hat{P}_{C_7}$
splits into
\[
  \mathfrak{F}_{\mathrm{struct}}
  = \mathbb{D}_{\mathrm{core}} + \mathbb{D}_{\mathrm{defect}},
\]
where $\mathbb{D}_{\mathrm{core}}$ contains autonomous admissible
transport and resonance, and $\mathbb{D}_{\mathrm{defect}}$ collects
closure mismatch, compensator activity, and off-sector suppression.

\paragraph{Tool (Exact-closure normal form).}
If $\mathfrak{R}_C = 0$, $[\mathcal{M}_{16},\hat{P}_{C_7}] = 0$, and
$\mathcal{H}_C = 0$ for the minimal compensator ansatz, then
$\mathfrak{F}_{\mathrm{struct}} = \mathbb{D}_{\mathrm{core}}$ and
$\mathbb{D}_{\mathrm{defect}} = 0$.
The visible compensator field vanishes in this regime; the kernel
$C = \Pi^{(\mathcal{M})}_{\mathrm{spec}}$
(Theorem~\texttt{thm:hc-spectral-projector} in the active-core PDF) remains active
(Remark~\texttt{rem:hc-shadow-vs-core} in the active-core PDF).

\paragraph{Tool (Mass-residue test).}
A proposed mass contribution is classified according to whether it
arises from: (i) a closed-core scalar reading of the projected resonance
operator; (ii) a weak closure-defect correction; (iii) a genuine
shifted-clock or projector residue
$\mathfrak{m}_{\mathrm{comm}}(\psi)
= \frac{\hbar}{c}|\langle\psi|[\hat{\mathcal{R}}_\beta,\hat{P}_{C_7}]|\psi\rangle|$.

\subsection*{Projection and regime tools}

\paragraph{Tool (Regime-strength test).}
A reduction of the structural master equation is: \emph{strong} if
exact closure plus direct slot restriction suffices; \emph{intermediate}
if projected algebra closure is additionally needed; \emph{effective} if
bilinear reconstruction, coarse-graining, or emergent parameter
generation are also required.
These levels correspond to the three steps of
Proposition~\texttt{prop:structural-lift-principle} in the active-core PDF.

\paragraph{Tool (Dirac gate).}
A Dirac-type reduction is available if: (i) exact closure holds;
(ii) the projected connection freezes locally to a trivial sector;
(iii) the projected resonance operator reduces to a scalar inertial
reading.

\paragraph{Tool (Lie-closure gate).}
A Yang--Mills-type reduction is available only if the projected
connection algebra closes on a finite-dimensional Lie algebra
$\mathfrak{g}$. Without that closure, the gauge interpretation is
structurally premature.

\paragraph{Tool (Einstein pipeline gate).}
An Einstein-type reduction requires all of: (i) exact closure on the
admissible sector; (ii) bilinear one-forms $E^a{}_\mu$ from projected
probe dynamics; (iii) maximal rank of the effective coframe; (iv)
metric-compatible geometric connection; (v) coarse-grained local
effective action with Einstein--Hilbert leading term. If any step is
missing, the Einstein block remains effective rather than
first-principles
(cf.\ Proposition~\texttt{prop:expected-probe-lensing-behaviour} in the active-core PDF).

\subsection*{Equivalence and anti-redundancy tools}

\paragraph{Tool (Stone--von Neumann compression rule).}
If two sectors are irreducible unitary realizations of the same
Weyl-type algebra on separable Hilbert spaces, differing only by
representation change, symmetry implementation, or frame reading, then
one representative suffices
(Theorem~\texttt{thm:unitary-equiv} in the active-core PDF).

\paragraph{Tool (Stone--von Neumann scope test).}
Do \emph{not} apply Stone--von Neumann as a shortcut if the passage
between two sectors involves: (i) controlled projection discarding
slots; (ii) coarse-graining over internal modes; (iii) bilinear
geometric reconstruction; (iv) emergent effective constants; (v) loss
of irreducibility. Stone--von Neumann compresses representation
multiplicity, not regime multiplicity.

\paragraph{Tool (Two-layer equivalence test).}
For candidate sectors $S_1, S_2$: (1)~\emph{Representation layer} ---
test whether they are unitary realizations of the same irreducible
Weyl-type kernel; if yes, SvN-compressible. (2)~\emph{Regime layer} ---
test whether one sector is obtained from the other by projection,
coarse-graining, geometric readout, or effective reconstruction; if yes,
additional structural work is required beyond Stone--von Neumann
(Definition~\texttt{def:readout-relative-lemma} in the active-core PDF).

\subsection*{Geometry-readout tools}

\paragraph{Tool (Coframe extraction gate).}
Given a closed triadic probe field $\Psi_\triangle$, define
$E^a{}_\mu := \operatorname{Re}(\bar\Psi_\triangle\,\Gamma^a\,
\mathcal{D}_\mu\Psi_\triangle)$.
This defines an effective coframe only if the resulting $4\times 4$
matrix has maximal rank.

\paragraph{Tool (Closure-to-geometry cleanliness test).}
If exact closure holds, probe bilinears are computed entirely inside
the invariant $C_7$-sector, so reconstructed metric data is not
contaminated by leakage terms. This is the minimal cleanliness
criterion for any geometric readout.

\paragraph{Tool (Geometry-reconstruction warning).}
Closure by itself does not yet produce Einstein geometry. It provides:
(i) an autonomous admissible sector; (ii) a well-defined projected
connection; (iii) a stable basis for bilinear readout. It does
\emph{not} yet provide automatically: (i) nondegeneracy of the
coframe; (ii) derivation of the effective constants; (iii) explicit
identification of the collective closure fields; (iv) justified
coarse-graining to a local geometric action. The bidirectional probe
lensing hypothesis (Section~\texttt{sec:bidirectional-probe-lensing} in the active-core PDF)
proposes a two-step route that may reduce this gap.

% Appendix mode is already active above; the following block continues the appendix material without re-entering appendix mode.

